problem:  
Let \((M,\mathcal{D},g)\) be a compact, equiregular **step‑2 contact sub‑Riemannian 5‑manifold** equipped with its Popp measure, and let \(\Delta\) be the associated sub‑Laplacian (symmetric with respect to the Popp measure). Denote by \(p_t(x,y)\) the corresponding heat kernel and assume that, as \(t\to 0^+\),  
\[
p_t(x,x)\sim \frac{1}{t^{Q/2}}\Bigl(1+\beta_1(x)t+\beta_2(x)t^2+\cdots\Bigr),\quad Q=6.
\]  
In the 3‑dimensional contact case, \(\beta_1,\beta_2\) can be explicitly expressed via the Tanaka–Webster curvature and torsion invariants; in dimension 5, these coefficients depend on additional curvature‑type tensors arising from the Levi form and the sub‑Riemannian connection.

**Open problem:**  
Assume that \((M,\mathcal{D},g)\) has tangent cone at every point isometric to the Heisenberg group \(\mathbb{H}^2\), and that the first three heat coefficients of the diagonal asymptotic match those of the flat model,  
\[
\beta_1(x)\equiv 0,\qquad \beta_2(x)\equiv 0,\qquad \beta_3(x)\equiv 0.
\]  
Is the manifold then necessarily locally flat, i.e. locally sub‑Riemannian isometric to \(\mathbb{H}^2\)?  
Equivalently, do the first three diagonal coefficients of the heat kernel uniquely determine local flatness in the 5D contact, step‑2 sub‑Riemannian category?

Why is it a "good" problem:  
This refinement directly targets a concrete and unsolved setting where partial asymptotic expansions are known but not fully classified. Restricting to five‑dimensional **contact** structures ensures that curvature tensors are explicitly defined (through Tanaka–Webster connections and higher‑order torsion), while the choice of \(k=3\) isolates the first as‑yet‑uncomputed layer of heat invariants beyond what is understood in 3D.  

The problem asks whether *finite spectral data*—specifically, the first three heat coefficients—suffice to detect local flatness when the nilpotent model is fixed as \(\mathbb{H}^2\). This question tightly links **analysis** (heat kernel asymptotics), **geometry** (sub‑Riemannian curvature invariants), and **rigidity theory** (local isometric classification). Unlike the general “can one hear the shape” type question, this version is concrete, focused on a specific dimension and structure, and sits precisely at the boundary between known and unknown results.  

A resolution would have deep implications: a positive answer would establish *finite‑order spectral rigidity* for 5D contact geometries, while a negative answer would reveal the existence of *spectrally invisible higher‑curvature effects*. Either outcome would open new ground in the ongoing development of sub‑Riemannian spectral geometry.